\section{Discussion and Conclusion}
\label{sec:conclusion}

We have introduced HospitAI, which brings together three analytical stages in hospitals based on data analysis: clustering patients using mixed-type clustering, sequence mining, and AI interpretation. The clustering produces a small set of understandable clusters of patients. Sequence mining, enhanced by geometric mean lift scoring and diverse filtering, provides clinically relevant paths rather than statistically frequent ones. Language generation converts numerical representations into short summaries and allows natural language querying.

There are also several limitations. The data is artificial, meaning the trends represent modeling in Synthea and do not depict an actual hospital's processes. We also have not conducted any form of clinical validation with medical practitioners, and the mined trends were evaluated based on qualitative evaluation rather than quantitatively compared to other miners with a benchmark for the same data set. The system depends on one external language model, which introduces challenges of expense and repeatability and poses a problem of governance that needs to be resolved before deployment at a hospital. The filter we used is tuned for the data, implying that some of the cleaning process is unique to this dataset.

Several next steps follow. First, we need to validate our approach on a real de-identified dataset and have clinicians evaluate our mined pathways. Second, a better evaluation of the sequence patterns would be beneficial, for instance by comparison with existing clinical guidelines or against other mining algorithms. Finally, a detailed analysis of the language layer is needed, particularly the quality of the SQL agent and of the produced summaries. Replacing the dataset-specific filtering with a more generic one would also make the pipeline more transferable. With appropriate preprocessing and postprocessing, off-the-shelf algorithms already give a hospital an affordable platform built on data it should already have.
